\documentclass[11pt]{article}
\usepackage[margin=1in]{geometry}
\usepackage{cmap}
\usepackage[T1]{fontenc}
\input{glyphtounicode}
\pdfgentounicode=1
\usepackage{booktabs,hyperref,amsmath}
\hypersetup{colorlinks=true,linkcolor=blue,urlcolor=blue}
\title{\textbf{Response to the fix list}}
\author{}\date{11 October 2026}
\sloppy
\begin{document}\maketitle\vspace{-3em}

Four of the five headline items are applied. All three checkers green afterwards: 63
headline quantities, 351 printed values, 25 superseded literals. 53 pages.

\section*{Two things in the fix list are wrong, and one of my passages was wrong again}

\textbf{The Fe$_3$C figure is mis-attributed.} The list credits the nearly 30\%
over-magnetisation of Fe$_3$C to Mej\'ia-Rodr\'iguez and Trickey. It is Fu and Singh's
result, which Mej\'ia-Rodr\'iguez and Trickey quote: ``In Fe--3C, \emph{they} found SCAN
gives nearly 30\% over-magnetization per three-iron-atom formula unit.'' The paper now
attributes it to Fu and Singh. Mej\'ia-Rodr\'iguez and Trickey's own contribution is
diagnostic --- tracing the error to the iso-orbital indicator --- and they are cited for
the $\ge$20\% bcc Fe figure and that diagnosis.

\textbf{The ``Relative'' column cannot be the better primary outcome.} The list suggests
checking it, and it was the right thing to check. It is documented as the leave-one-out
hull, and it is \emph{empty in both released files}: zero populated rows of 554{,}054 in
the summary file and zero of 139{,}541 in the recomputation file. The paper now says so in
one sentence, and \texttt{phase40} asserts it, so a future run notices if DeepMind ever
populates it.

\textbf{My \S4.13 passage was wrong a second time.} Draft 24 said recovering the
Materials Project sub-hull margin ``requires the competing entries of each chemical system
and not the four summary fields we retrieved.'' That is false. The thermo documents expose
\texttt{decomposition\_enthalpy} --- the phase-separation energy of Bartel et al., computed
by removing each entry from its own phase diagram, negative for stable entries and
populated for all of them --- and
\texttt{equilibrium\_rea}\-\texttt{ction\_energy\_per\_atom}, the same margin restricted to
entries flagged stable. One extra query against the same
endpoint supplies the covariate. The passage now says the omission is mine and is cheap to
repair. That is twice I have written a wrong account of this one gap.

\section*{P0-1. The reference sets are documented}

Verified against \texttt{DATASET.md}: ``All'' is the distance to the convex hull of all
entries, ``Relative'' the hull of all entries except the current one, ``MP'' the Materials
Project hull, ``MP OQMD'' Materials Project plus OQMD including recalculations. \S3.1 now
states all four with a citation, and \S4.16's three paragraphs of hedging are gone ---
``unidentified alternative reference set'', ``our reading, which we state as an inference'',
``the question is settled by asking the dataset authors'' all deleted. It reads as a
sensitivity analysis against a known denser hull, and the intermetallic rise from 9.2\% to
22.3\% at class level and 4.0\% to 14.5\% in the headline cell is now stated as evidence
that the MP-referenced hull is least complete for intermetallics.

Two caveats survive because they are true. \texttt{DATASET.md} documents the
\emph{summary} file and does not define the recomputation file's columns, which carry the
same four names and are where the 25.8\% comes from; the paper says it takes the documented
meanings as applying to both, on the strength of the identical names. And a denser hull
still does not identify which phases were added to which chemistry.

\section*{P0-2. The r$^2$SCAN protocol}

New paragraph in \S3.1: three jobs per material (PBE pre-relaxation with PBE54 potentials,
r$^2$SCAN relaxation, r$^2$SCAN static), \texttt{MPScanRelaxSet} and
\texttt{MPScanStaticSet} settings, structures re-relaxed rather than evaluated at the PBE
geometry, PBE52/PBE54 PAW datasets that can differ from those behind the PBE labels, and
ferromagnetic initialisation throughout. The sentence that matters: a verdict change
reflects functional, correction scheme, PAW dataset, relaxed geometry and possibly
converged magnetic state together. The PAW-dataset confound was absent from the paper
entirely.

\section*{P0-3. Licence}

Verified at source: GNoME data under CC BY-NC 4.0, code under Apache 2.0. A licensing
paragraph is now in the data statement and \texttt{results/PROVENANCE.md} is corrected ---
it previously said only that the GNoME inputs were governed by ``whatever terms DeepMind
applies to them''.

The redistribution check passes. No file in \texttt{results/} carries GNoME per-material
rows; the largest GNoME-derived result is 83 rows of per-element summaries. The two large
per-material files are Materials Project extracts under CC BY 4.0, which is what
\texttt{PROVENANCE.md} already documented. The non-commercial clause against a company
affiliation is flagged in both places and left for the author; it is not a question this
document can answer.

\section*{P1-1. The magnetism literature}

All four references verified, with full author lists and the attributed findings checked
against the papers. ``The data do not support that reading'' is gone. The paragraph now
reads: an element-presence indicator cannot separate $U$-list membership from 3$d$
magnetism, and one magnetic mechanism is well documented --- Ekholm et al.\ on bcc Fe, hcp
Co and fcc Ni with SCAN worse than PBE for Ni and Co; Fu and Singh on V, Co, Ni, Pd and
Fe$_3$C; Mej\'ia-Rodr\'iguez and Trickey on $\ge$20\% in bcc Fe; Tran et al.\ on
intermetallic ferromagnets; Kothakonda et al.\ on r$^2$SCAN exceeding SCAN, iron
overestimated by 24\% against 17\%.

Three features of the results are stated as consistent with it: the Fe, Co, Cr, Mn ordering
of the per-element rates, the concentration of the iron and cobalt excess in metallic anion
classes, and its near-absence in oxides. Nickel is no longer anomalous --- its nearly filled
majority band limits further enhancement. Molybdenum and tungsten remain unexplained and
are still named as such. The paragraph ends by saying the mechanism is untested, that the
released files carry no magnetisation field, and what the discriminating calculation would
be.

\section*{P2. Editing defects, and the text layer}

Applied: the reference-set sentence is in the abstract; both figure titles corrected inside
the images and regenerated (``red = GNoME flips more, green = GNoME flips less'';
``Flip rate by PBE band-gap bin and 3$d$ content''); ``nearly eightfold range'' replaced by
the 26.7-point spread in \S4.1 as well as the abstract; the per-element heading now says
``against the MP reference''; every remaining ``registered'' is ``pre-committed'', including
the \S3.8 heading; the unmatched-identifier sentence now says 2.0 points lower and not
negligible at $n=22{,}498$; the ``The time of writing'' typo fixed; the 84\% and 86.8\%
comparison now says the source gives no denominators or reference hull.

\textbf{The ligature problem was real and is fixed.} \texttt{lmodern} is not installed in
this build environment, so the prescribed preamble fails; \texttt{cmap} plus
\texttt{glyphtounicode} does the job. Measured on the built PDF: ``flip rate'' 35
occurrences, ``difference'' 28, ``workflow'' 19, ``coefficient'' 10, and zero occurrences of
the broken forms the old text layer produced. On Draft 24 the same test found zero of the intact forms.

\section*{Queued}

P1-3, pymatgen version and the no-corrections-recomputed statement. P1-4, lanthanide PAW
labels from \texttt{MPRelaxSet.yaml} --- concrete and not yet done. P1-5, the prior-art
search. P1-2's actual re-pull, which needs the machine that holds the credentials; the
Materials Project API is unreachable from here.

The cut plan to about 14 pages plus supplement is the largest remaining item and is
deliberately last: this pass added a protocol paragraph, a mechanism paragraph, a licensing
paragraph and four references. The supplement has to exist before the cut can land, and the
cut is where the 53 pages come down.

From P3: fixing the propensity model with the spline already in the paper and trimming to
the overlap region, BCa intervals for the two headline cells and the Fe oxide excess,
reliability curve and Brier score on held-out clusters, and the finite-cluster correction in
the sandwich estimator. All four are feasible from cached data. The magnetism calculation
and the third reference hull are not.

\end{document}
